\frontmatterchapter{Acknowledgements}

Sincere thanks go to the supervisor, Luís Miguel Moreira Lino Ferreira, PhD, and the co-supervisor, Armando Jorge Sousa, PhD, for their guidance and trust. Throughout many meetings, they made time to listen to my ideas, discuss different possibilities, and return to questions as the research developed. These conversations helped shape the work and were among the most meaningful parts of the doctoral experience.

Particular thanks are due to the industrial partners CleanWatts, especially Eng.~Jorge Landeck, and i-Charging, especially Eng.~Paulo Rodrigues. Their experience with energy systems and electric-vehicle charging brought a practical understanding of the constraints encountered during the work in OPEVA.

The students whose theses and research projects developed alongside this work made supervision a source of learning. Their questions opened up different ways of approaching related problems, and the opportunity to accompany their projects from an initial idea through development and validation was especially valuable.

\draftnote{Add the names of the students the author wishes to thank individually.}

Friends and colleagues at FEUP made the day-to-day experience of the doctorate easier to share. João Vitorino and Tiago Dias deserve a special mention for the conversations, shared frustrations, and laughter that accompanied these years.

My parents, Vera and Abílio, deserve a particularly warm word of thanks. They provided the stability and opportunities on which this path was built, long before the doctorate began. Their encouragement sustained choices that sometimes demanded patience and sacrifice from them, and their example taught the value of work, responsibility, and perseverance. The affection and presence of my grandparents were equally important, giving a sense of family that extended far beyond the years devoted to this thesis.

The final and most personal thanks go to Inês. She shared the achievements of the doctorate and lived closely with the pressure and exhaustion it brought into everyday life. There were periods when the work took up too much space, and its frustrations followed it home as impatience, distraction, or withdrawal. Her support continued through those periods, even when it was difficult to give. She made time to listen and helped restore perspective when a research problem seemed to occupy everything. She celebrated progress that was too easily dismissed, offered encouragement when confidence faltered, and knew when a challenge would help more than reassurance. Her belief remained steady through long periods of doubt and slow progress. The patience and affection behind that support meant a great deal to the completion of this work and to the shared life that continued around it. The order of a finished thesis conceals much of the uncertainty and repetition involved in writing it. Inês accompanied that less orderly part of the work too, with a patience and affection for which these pages offer only a small acknowledgement.

Nes, Pai, Mae, Aida, Vo Ek e Vo Fonso.
